\section*{الفصل \ref{ch:g7:fractions} --- الكسور: المقارنة والجمع}

\begin{solution}{exo:g7:fractions:1}
$\dfrac{2}{5} = \dfrac{8}{20}$; \quad
$\dfrac{18}{24} = \dfrac{3}{4}$; \quad
$\dfrac{7}{3} = \dfrac{28}{12}$; \quad
$\dfrac{5}{9} = \dfrac{20}{36}$.
\end{solution}

\begin{solution}{exo:g7:fractions:2}
$\dfrac{12}{16} = \dfrac34$; \quad
$\dfrac{30}{45} = \dfrac23$; \quad
$\dfrac{27}{36} = \dfrac34$; \quad
$\dfrac{48}{12} = 4$.
\end{solution}

\begin{solution}{exo:g7:fractions:3}
$\dfrac{15}{35} = \dfrac37$ و $\dfrac{21}{49} = \dfrac37$: متساويان.

$\dfrac{16}{28} = \dfrac47$ و $\dfrac{20}{36} = \dfrac59$: وللمقارنة،
$\frac47 = \frac{36}{63}$ و $\frac59 = \frac{35}{63}$ ---
غير متساويين.
\end{solution}

\begin{solution}{exo:g7:fractions:4}
$\dfrac{7}{11} < \dfrac{9}{11}$ (\omterm{def:g4:fractions:def}{المقام} نفسه).

$\dfrac56 = \dfrac{15}{18} > \dfrac{11}{18}$.

$\dfrac{13}{12} > 1$ بينما $\dfrac{19}{20} < 1$، إذن
$\dfrac{13}{12} > \dfrac{19}{20}$.
\end{solution}

\begin{solution}{exo:g7:fractions:5}
$\dfrac59 + \dfrac29 = \dfrac79$; \quad
$\dfrac{11}{7} - \dfrac47 = \dfrac77 = 1$; \quad
$\dfrac58 + \dfrac78 = \dfrac{12}{8} = \dfrac32$.
\end{solution}

\begin{solution}{exo:g7:fractions:6}
$\dfrac23 + \dfrac56 = \dfrac46 + \dfrac56 = \dfrac96 = \dfrac32$.

$\dfrac{7}{10} - \dfrac25 = \dfrac{7}{10} - \dfrac{4}{10} =
\dfrac{3}{10}$.

$\dfrac34 + \dfrac{5}{12} = \dfrac{9}{12} + \dfrac{5}{12} =
\dfrac{14}{12} = \dfrac76$.
\end{solution}

\begin{solution}{exo:g7:fractions:7}
$3 - \dfrac54 = \dfrac{12}{4} - \dfrac54 = \dfrac74$; \quad
$1 + \dfrac38 = \dfrac{11}{8}$; \quad
$2 - \dfrac76 = \dfrac{12}{6} - \dfrac76 = \dfrac56$.
\end{solution}

\begin{solution}{exo:g7:fractions:8}
$5 \times \dfrac{3}{10} = \dfrac{15}{10} = \dfrac32$; \quad
$\dfrac78 \times 4 = \dfrac{28}{8} = \dfrac72$; \quad
$\dfrac23$ من $45$ هو $\dfrac{2 \times 45}{3} = \dfrac{90}{3} = 30$.
\end{solution}

\begin{solution}{exo:g7:fractions:9}
معًا: $\dfrac14 + \dfrac38 = \dfrac28 + \dfrac38 = \dfrac58$ من
البيتزا. \omterm{def:g3:division:remainder}{والباقي}: $1 - \dfrac58 = \dfrac38$.
\end{solution}

\begin{solution}{exo:g7:fractions:10}
المسكوب: $2 \times \dfrac16 = \dfrac26 = \dfrac13$ لتر. \omterm{def:g3:division:remainder}{والباقي}:
\[
\frac34 - \frac13 = \frac{9}{12} - \frac{4}{12} = \frac{5}{12}
\text{ لتر}.
\]
\end{solution}

\begin{solution}{exo:g7:fractions:11}
جمع المقدار الموجب $\frac12$ إلى $\frac35$ يجب أن يعطي نتيجة
\emph{أكبر} من $\frac35$. لكن $\frac47 < \frac35$ (فعلًا
$\frac47 = \frac{20}{35}$ و $\frac35 = \frac{21}{35}$): \omterm{def:g1:addition:def}{فالمجموع}
المدّعى أصغر من أحد حدّيه --- وهذا مستحيل.
\end{solution}

\begin{solution}{exo:g7:fractions:12}
خطوة خطوة:
\[
\frac12 + \frac14 = \frac34, \qquad
\frac34 + \frac18 = \frac68 + \frac18 = \frac78, \qquad
\frac78 + \frac1{16} = \frac{14}{16} + \frac1{16} = \frac{15}{16}.
\]
والمجاميع الجزئية $\frac12, \frac34, \frac78, \frac{15}{16}, \dots$
تقترب في كل مرة نصف المسافة الباقية إلى $1$: ففي كل خطوة يُقطع الجزء
الناقص إلى نصفين. وتقترب المجاميع من $1$ دون أن تبلغه أبدًا.
\end{solution}

\begin{solution}{pb:g7:fractions:1}
\textbf{1.} (أ) $\frac12 + \frac14 + \frac14 = \frac24 + \frac14
+ \frac14 = \frac44 = 1$: مشروع. (ب) $\frac14 + \frac18 +
\frac18 + \frac12 = \frac28 + \frac18 + \frac18 + \frac48 =
\frac88 = 1$: مشروع.

\textbf{2.} $\frac12 + \frac14 + \frac18 = \frac48 + \frac28 +
\frac18 = \frac78$: فتنقص علامة ثمن واحدة.

\textbf{3.} الميزان هو $1 = \frac{16}{16}$: أي ست عشرة علامة
جزء. والثمن هو $\frac18 = \frac{2}{16}$: أي جزأين.

\textbf{4.} النصف المنقوط: $\frac12 + \frac14 = \frac34$. والربع
المنقوط: $\frac14 + \frac18 = \frac28 + \frac18 = \frac38$.

\textbf{5.} النصف المنقوط يساوي $\frac34$ (السؤال 4):
أي ميزان فالس واحد بالضبط. وملء آخر، مثلًا: ربع
$+$ ربع $+$ ربع ($\frac14 + \frac14 + \frac14 =
\frac34$)، أو ربع $+$ ثمن $+$ ثمن $+$ ربع
($\frac14 + \frac18 + \frac18 + \frac14 = \frac{2 + 1 + 1 +
2}{8} = \frac68 = \frac34$).

\textbf{6.} $\frac38 + \frac18 + \frac14 = \frac38 + \frac18 +
\frac28 = \frac68 = \frac34$: فتنقص علامة ربع واحدة
($\frac14$).

\textbf{7.} الثمن المنقوط: $\frac18 + \frac{1}{16} =
\frac{2}{16} + \frac{1}{16} = \frac{3}{16}$. والربع:
$\frac{4}{16}$. فالربع أطول.

\textbf{8.} $\frac12 + \frac18 = \frac48 + \frac18 = \frac58$،
أي $\frac{10}{16}$: أي عشرة أجزاء من ستة عشر.

\textbf{9.} $1 - \frac{11}{16} = \frac{16}{16} - \frac{11}{16}
= \frac{5}{16}$ من الميزان صمت.

\textbf{10.} المثالان كلاهما مجموعه $\frac{2 + 1 + 1}{16} =
\frac{4}{16}$: مشروع. وكل الإيقاعات التي تملأ أربعة أجزاء
بقطع من $2$ و $1$:
\[
2{+}2, \quad 2{+}1{+}1, \quad 1{+}2{+}1, \quad 1{+}1{+}2, \quad
1{+}1{+}1{+}1 :
\]
أي خمسة إيقاعات.

\textbf{11.} الإيقاع الذي يملأ $n$ جزءًا ينتهي إما بجزء
واحد --- وما يسبقه يملأ $n - 1$ --- وإما
بثمن --- وما يسبقه يملأ $n - 2$. وكل إيقاع
يُعدّ مرة واحدة بالضبط بهذه الطريقة، إذن
$\text{عدد الطرق}(n) = \text{عدد الطرق}(n-1) + \text{عدد الطرق}(n-2)$. والجدول:
\begin{center}
\renewcommand{\arraystretch}{1.2}
\begin{tabular}{c|cccccccc}
$n$ & $1$ & $2$ & $3$ & $4$ & $5$ & $6$ & $7$ & $8$ \\
\hline
عدد الطرق & $1$ & $2$ & $3$ & $5$ & $8$ & $13$ & $21$ & $34$ \\
\end{tabular}
\end{center}
ويمكن قرع علامة النصف ($8$ أجزاء) في $34$ طريقة.
(والحالة $n = 4$ تستعيد الإيقاعات الخمسة في السؤال 10.)

\textbf{12.} كل عدد هو \omterm{def:g1:addition:def}{مجموع} العددين اللذين قبله:
$3 + 5 = 8$، و $5 + 8 = 13$، و $8 + 13 = 21$، و $13 + 21 = 34$ ---
وهي القاعدة التي برهن عليها السؤال 11، عمرها ألف سنة وتُعرف الآن
بقاعدة فيبوناتشي.

\textbf{13.} $\frac16 + \frac{1}{12} = \frac{2}{12} +
\frac{1}{12} = \frac{3}{12} = \frac14$: أي المدّة نفسها التي
لثمنين \omterm{def:g6:lines:objects}{مستقيمين} ($\frac18 + \frac18 = \frac14$). فتبديل
التأرجح مشروع.

\textbf{14.} كل علامة ثلاثية تدوم $\frac{1}{12}$: والتحقّق،
$3 \times \frac{1}{12} = \frac{3}{12} = \frac14$
(\cref{prop:g7:fractions:mult}). ومن أجل علامة نصف: تدوم كل واحدة
$\frac16$، لأن $3 \times \frac16 = \frac36 = \frac12$.

\textbf{15.} ربع منقوط $+$ ربع $+$ ربع: بالأثمان،
$3 + 2 + 2 = 7$ أثمان $= \frac78$: مشروع. والأنصاف والأرباع
تساوي $4$ و $2$ من الأثمان --- وكلاهما عدد \emph{زوجي}. \omterm{def:g1:addition:def}{ومجموع}
\omterm{def:g3:numbers:evenodd}{الأعداد الزوجية} زوجي، لكن ميزان سبعة أثمان يحتاج
$7$ أثمان، وهو \omterm{def:g3:numbers:evenodd}{عدد فردي}: فهذا مستحيل. ولملء ميزان فردي،
لا بدّ أن يظهر شيء فردي --- ثمن أو علامة منقوطة.
\end{solution}
